\section{The local WSPD death-time implementation}\label{sec:wspd}

This appendix discharges the one implementation step deferred by \cref{sec:dt}: realizing a locally
navigable Euclidean $(1+\rho)$-spanner of a point set $X\subseteq[\Delta]^2$ under the range-counting
oracle, so that the clipped death-time sampler of \cref{sec:dt-sampler} can be run from a uniformly
sampled vertex at a cost we control. We follow the well-separated pair decomposition (WSPD) construction
of Driemel et al.~\cite{DMOSW25} (whose spanner facts are from Har-Peled~\cite{HarPeled11} and whose
sampling primitive is from Monemizadeh~\cite{Monemizadeh23}), and make every step concrete as a sequence
of range counts. We use the term \emph{cell} for a node of the quadtree $\mathcal T$ on $[\Delta]^2$; a
cell of side $r$ has radius $r(c)\le\sqrt2\,r$, and $d(c,c')$ is the distance between cell centers. All
counts below are on the input $P$; when $X=Q$ is a snapped support, a $Q$-cell is nonempty exactly when
the aligned $P$-region is nonempty, so each emptiness test is still one count on $P$ (\cref{sec:dt}).

\subsection{Uniform point sampling under the oracle}\label{sec:prelim-sample}

We first record the sampling primitive used here and in \cref{sec:support} (referenced from
\cref{sec:support-oracle}). It is \emph{telescoping sampling}~\cite{Monemizadeh23}.

\begin{lemma}[Uniform sampling]\label{lem:telescoping}
Given a query rectangle $R$ that is a union of quadtree cells with $\rcount{R}>0$, one can return a point
of $P\cap R$ drawn uniformly at random using $O(\log\Delta)$ range-counting queries.
\end{lemma}

\begin{proof}
Walk down $\mathcal T$ from the cell(s) covering $R$. At a cell $c$ with children $c_1,\dots,c_4$, query
the four counts $\rcount{c_i}$ (which sum to $\rcount c$) and descend into child $c_i$ with probability
$\rcount{c_i}/\rcount c$. After $O(\log\Delta)$ levels the walk reaches a single grid point, which by the
chain of conditional probabilities is uniform on $P\cap R$. Each level spends $O(1)$ counts, so the total
is $O(\log\Delta)$.
\end{proof}

Two consequences are used repeatedly. First, a uniform vertex of $X$ (a death-time seed) costs
$O(\log\Delta)$ counts. Second, the \emph{representative} of a cell, which we now define, is computable in
$O(\log\Delta)$ counts.

\subsection{The spanner realized by range counts}\label{sec:wspd-spanner}

\paragraph{Representatives.}
For a nonempty cell $c$, its representative $\rep(c)$ is the point of $P\cap c$ that is smallest in the
$z$-order (the order in which a quadtree traversal visits leaves). It is found by descending $\mathcal T$
from $c$, at each level entering the $z$-first nonempty child; one emptiness count per child decides this,
so $\rep(c)$ costs $O(\log\Delta)$ counts. The $z$-order is a total order on $[\Delta]^2$, so $\rep(c)$ is
unique and the same cell always returns the same representative across every search; this fixes all ties
and makes the spanner below a single well-defined graph.

\paragraph{Well-separated pairs.}
A pair of cells $(c,c')$ is \emph{well-separated} at parameter $\rho$ if $\rho\cdot d(c,c')\ge
\max\{r(c),r(c')\}$. The WSPD $\mathcal W$ of $X$ is the family produced by the Callahan--Kosaraju
recursion~\cite{CK95} on $\mathcal T$, of size $O(\rho^{-2}n)$~\cite{DMOSW25}. We never materialize
$\mathcal W$. Two facts let us work with it locally:
\begin{enumerate}
\item Membership $(c,c')\in\mathcal W$ is decided in $O(1)$ counts: well-separation is a geometric test
on the two cells (no query), and the recursion places $(c,c')$ in $\mathcal W$ exactly when $(c,c')$ is
well-separated, both cells are nonempty, and the parent pair is not well-separated; the three emptiness
tests are $O(1)$ counts.
\item (\cite{HarPeled11}, restated in \cite{DMOSW25}.) For every $(c,c')\in\mathcal W$, $r(c')/2\le
r(c)\le 2r(c')$, and every point of $[\Delta]^2$ lies in $O(\rho^{-2}\log\Delta)$ pairs of $\mathcal W$.
\end{enumerate}
The \emph{spanner} $S=(X,E(S))$ adds, for each $(c,c')\in\mathcal W$, the edge $\{\rep(c),\rep(c')\}$ of
length $d(c,c')$. By the second fact every vertex has $O(\rho^{-2}\log\Delta)$ incident spanner edges, so
$S$ is sparse.

\paragraph{Redundant-pair suppression.}
The Callahan--Kosaraju recursion run on the quadtree of $[\Delta]^2$ (not on the points) can emit two
pairs $(c,c')$ and $(\bar c,\bar c')$ with $P\cap c=P\cap\bar c$ and $P\cap c'=P\cap\bar c'$: geometrically
distinct cells carrying the same point sets. We suppress all but the pair whose two cells are smallest in
$z$-order among those with equal point sets. The test ``$P\cap c=P\cap\bar c$'' for two candidate cells
of comparable side reduces, by \cref{lem:telescoping}'s walk, to comparing $\rep$ and the count
$\rcount c$, and adds $O(\log\Delta)$ counts per pair touched. Suppression makes each point-pair $\{x,y\}$
covered by a unique surviving pair, so the surviving family still has size $O(\rho^{-2}n)$ and the
per-point bound $O(\rho^{-2}\log\Delta)$ of fact~(2) holds for the spanner $S$.

\paragraph{Stretch.}
The spanner inherits the standard guarantee.

\begin{lemma}[Spanner stretch \cite{HarPeled11,DMOSW25}]\label{lem:stretch}
$\wt(\MST(S))\le(1+\rho)\,\wt(\MST(X))$, and for every threshold $t$,
$c_X(t)\le c_S(t)\le c_X(t/(1+\rho))$.
\end{lemma}

The first inequality is Lemma~23 of \cite{DMOSW25}. The second is the threshold-graph form: an $S$-edge of
length $d(c,c')$ joins two points at distance $\le(1+\rho)\,d(c,c')$ and $\ge d(c,c')-r(c)-r(c')\ge
d(c,c')/(1+\rho)$ by well-separation, so an $S$-edge of length $\le t$ certifies an $X$-pair within
$(1+\rho)t$, and an $X$-pair within $t/(1+\rho)$ is spanned by an $S$-edge of length $\le t$. Integrating
the second inequality against \eqref{eq:cluster-integral} gives \eqref{eq:spanner-stretch}.

\subsection{Local navigability}\label{sec:wspd-nav}

The death-time search needs, from a vertex of $S$, its spanner neighbors within a given length. We provide
this at each dyadic length scale.

\begin{lemma}[Neighbor oracle]\label{lem:navigate}
Let $c$ be a cell of side $r/2$. All cells $c'$ of side $r/2$ such that $S$ has an edge of length at most
$r$ between a point of $c$ and a point of $c'$ can be found using $O(\rho^{-2})$ range-counting queries.
\end{lemma}

\begin{proof}
This is Lemma~24 of \cite{DMOSW25}; we restate the realization. If an $S$-edge of length $\le r$ joins $c$
and $c'$ then $d(c,c')\le 3r$, and there are $O(1)$ cells $c'$ of side $r/2$ with $d(c,c')\le 3r$. For each
such candidate we must decide whether some surviving pair of $\mathcal W$ links a subcell of $c$ to a
subcell of $c'$ at distance $\le r$. A direct WSPD recursion on $(c,c')$ could descend $\Theta(n)$ times,
so we truncate at side $2^{i^*}$ with $2^{i^*}\le\rho r/10<2^{i^*+1}$: declare a \emph{witness} a pair
$(\bar c,\bar c')$ of nonempty cells of side $\ge 2^{i^*}$, with $\bar c\subseteq c$, $\bar c'\subseteq
c'$, satisfying either (i) $(\bar c,\bar c')\in\mathcal W$ and $d(\bar c,\bar c')\le r$, or (ii) $(\bar
c,\bar c')$ not well-separated and $d(\bar c,\bar c')+r(\bar c)+r(\bar c')\le r$. There are
$O(\rho^{-2})$ subcells of side $\ge 2^{i^*}$ inside $c$ (and inside $c'$), so $O(\rho^{-4})$ candidate
witness pairs; a modification of the recursion that tests these pairs by membership and emptiness
(each $O(1)$ counts) and prunes well-separated branches early completes the search in $O(\rho^{-2})$
counts. Correctness of the truncation (a witness exists iff $(c,c')$ is a yes-instance) is the argument of
\cite[Lemma~24]{DMOSW25}: a too-small defining pair forces its parent pair to be non-well-separated, which
brings the endpoints within $2r/5$ and so produces a side-$2^{i^*}$ witness.
\end{proof}

\subsection{The bottleneck death-time search}\label{sec:wspd-search}

We can now compute the clipped death time $X_L(v)=\min\{\tau(v),L\}$ of a seed $v$ (\cref{sec:dt-sampler}).
Assign each vertex an independent uniform rank in $[0,1]$, drawn lazily and stored by cell identity, so
ranks are consistent and a.s.\ distinct (ties, of probability zero, are broken by $z$-order). Recall
$\tau(v)$ is the bottleneck (min--max-edge) distance from $v$ to the root or to any lower-ranked vertex.

\paragraph{The search.}
Run a bottleneck Dijkstra from $v$ over the dyadic length scales $r=(1+\rho),(1+\rho)^2,\dots$ up to $L$.
Maintain the set $R$ of reached vertices, initialized to $\{v\}$. At scale $r$, for every reached vertex
$u$ apply \cref{lem:navigate} (at the cell side matching $u$) to enumerate the spanner neighbors of $u$
joined by an edge of length $\le r$, and add them to $R$. Halt and output $\tau(v)=r$ the first time $R$
contains a vertex of rank smaller than $\mathrm{rank}(v)$; if the scale passes $L$ first, output
$X_L(v)=L$. Because edges are processed in increasing length, the level at which a lower-ranked vertex
first enters $R$ is the bottleneck value, so the procedure returns $\min\{\tau(v),L\}$.

\paragraph{Per-search cost.}
The number of vertices extracted before a lower-ranked one appears is small in expectation: order the
vertices by bottleneck distance from $v$; the $j$-th is extracted only if $v$ has the smallest rank among
the first $j$ (else the search would have halted earlier), an event of probability $1/j$, so the expected
number of extractions is $\sum_{j\ge1}1/j$ truncated at $|V|$, i.e.\ $H_{|V|}=O(\log|V|)$ (the same
harmonic bound as the leader estimator, \cref{lem:leader}). Each extracted vertex is expanded over
$O(\rho^{-1}\log\Delta)$ scales, each costing $O(\rho^{-2})$ counts by \cref{lem:navigate}, so a search
costs $O(\rho^{-3}\log\Delta)\cdot O(\log|V|)=\Otil_\rho(1)$ range counts in expectation. To bound the
worst case we cap the extractions at $\Theta(\log|V|\cdot\log\tfrac1{\delta'})$; by Markov the cap is hit
with probability $\le\delta'$, and on that event we output $X_L(v)=L$. Since $X_L\le L$ always, capping
only lowers the reported value on a $\delta'$-fraction of seeds, an effect absorbed into the additive
budget of \cref{sec:main} by taking $\delta'$ a small constant times the target relative accuracy.

\subsection{The global query cap}\label{sec:wspd-cap}

The sampler of \cref{sec:dt-sampler} draws seeds uniformly from $V$ (\cref{lem:telescoping},
$O(\log\Delta)$ counts each) and runs the search above on each. For the bulk term $A_L(P)$ the assembly
uses $s_P=O(\xi^{-2}n^{1/3})$ seeds (\cref{sec:main}); at $\Otil_\rho(1)$ counts per search the bulk
sampler costs $\Otil_\xi(n^{1/3})$ counts. For $A_L(Q)$ the sampler uses $s_Q=O_\xi(1)$ support seeds,
each drawn and explored at the support-sampling cost $\Otil(n/K)=\Otil(n^{1/3})$ of
\cref{sec:support-oracle}, again $\Otil_\xi(n^{1/3})$. Redundant-pair suppression keeps every vertex
degree $O(\rho^{-2}\log\Delta)$, and the extraction cap bounds each search deterministically, so no search
runs away and the total over the whole algorithm is $\Otil_\eps(n^{1/3})$, matching the count claimed in
\cref{sec:dt} and \cref{sec:main}.
